\title{2019 年高考数学 全国I卷-理科}
\date{}
\maketitle

\section*{一、单项选择题}
本题共12小题，每小题5分，共60分。在每小题给出的四个选项中，只有一项是符合题目要求的。

\begin{question}[points=5]
已知集合 $M = \{x \mid -4 < x < 2\}$, $N = \{x \mid x^2 - x - 6 < 0\}$, 则 $M \cap N =$
\begin{choices}
  \item $\{x \mid -4 < x < 3\}$
  \item $\{x \mid -4 < x < -2\}$
  \item $\{x \mid -2 < x < 2\}$
  \item $\{x \mid 2 < x < 3\}$
\end{choices}
\begin{solution}
\textbf{答案：}$C$

\textbf{分析：}本题考查集合的交集和一元二次不等式的解法，渗透了数学运算素养．采取数轴法，利用数形结合的思想解题．

\textbf{详解：}由题意得，$M = \{x \mid -4 < x < 2\}$, $N = \{x \mid -2 < x < 3\}$, 则 $M \cap N = \{x \mid -2 < x < 2\}$．故选 $C$．
\end{solution}
\end{question}

\begin{question}[points=5]
设复数 $z$ 满足 $|z - i| = 1$, $z$ 在复平面内对应的点为 $(x, y)$, 则
\begin{choices}
  \item $(x + 1)^2 + y^2 = 1$
  \item $(x - 1)^2 + y^2 = 1$
  \item $x^2 + (y - 1)^2 = 1$
  \item $x^2 + (y + 1)^2 = 1$
\end{choices}
\begin{solution}
\textbf{答案：}$C$

\textbf{分析：}本题考点为复数的运算，为基础题目，难度偏易．此题可采用几何法，根据点 $(x, y)$ 和点 $(0, 1)$ 之间的距离为 $1$，可选正确答案 $C$．

\textbf{详解：}$z = x + yi$, $z - i = x + (y - 1)i$, $|z - i| = \sqrt{x^2 + (y - 1)^2} = 1$, 则 $x^2 + (y - 1)^2 = 1$．故选 $C$．
\end{solution}
\end{question}

\begin{question}[points=5]
已知 $a = \log_2 0.2$, $b = 2^{0.2}$, $c = 0.2^{0.3}$, 则
\begin{choices}
  \item $a < b < c$
  \item $a < c < b$
  \item $c < a < b$
  \item $b < c < a$
\end{choices}
\begin{solution}
\textbf{答案：}$B$

\textbf{分析：}运用中间量 $0$ 比较 $a$, $c$，运用中间量 $1$ 比较 $b$, $c$．

\textbf{详解：}$a = \log_2 0.2 < \log_2 1 = 0$, $b = 2^{0.2} > 2^0 = 1$, $0 < 0.2^{0.3} < 0.2^0 = 1$, 则 $0 < c < 1$, $a < c < b$．故选 $B$．
\end{solution}
\end{question}

\begin{question}[points=5]
古希腊时期，人们认为最美人体的头顶至肚脐的长度与肚脐至足底的长度之比是 $\frac{\sqrt{5} - 1}{2}$ ( $\frac{\sqrt{5} - 1}{2} \approx 0.618$ ，称为黄金分割比例)，著名的“断臂维纳斯”便是如此．此外，最美人体的头顶至咽喉的长度与咽喉至肚脐的长度之比也是 $\frac{\sqrt{5} - 1}{2}$ ．若某人满足上述两个黄金分割比例，且腿长为 $105$ cm，头顶至脖子下端的长度为 $26$ cm，则其身高可能是
\begin{choices}
  \item $165$ cm
  \item $175$ cm
  \item $185$ cm
  \item $190$ cm
\end{choices}
\begin{solution}
\textbf{答案：}$B$

\textbf{分析：}理解黄金分割比例的含义，应用比例式列方程求解．

\textbf{详解：}设人体脖子下端至腿根的长为 $x$ cm，肚脐至腿根的长为 $y$ cm，则 $\frac{26}{x} = \frac{26 + x}{y + 105} = \frac{\sqrt{5} - 1}{2}$，得 $x \approx 42.07$ cm, $y \approx 5.15$ cm．又其腿长为 $105$ cm，头顶至脖子下端的长度为 $26$ cm，所以其身高约为 $42.07 + 5.15 + 105 + 26 = 178.22$，接近 $175$ cm．故选 $B$．
\end{solution}
\end{question}

\begin{question}[points=5]
函数 $f(x) = \frac{\sin x + x}{\cos x + x^2}$ 在 $[-\pi, \pi]$ 的图像大致为

\begin{center}
  \includegraphics[width=.38\linewidth]{question-5-1.jpg}
\end{center}
\begin{choices}
  \item A． (图片)
  \item B． (图片)
  \item C． (图片)
  \item D． (图片)
\end{choices}
\begin{solution}
\textbf{答案：}$D$

\textbf{分析：}先判断函数的奇偶性，得 $f(x)$ 是奇函数，排除 $A$，再注意到选项的区别，利用特殊值得正确答案．

\textbf{详解：}由 $f(-x) = \frac{\sin(-x) + (-x)}{\cos(-x) + (-x)^2} = \frac{-\sin x - x}{\cos x + x^2} = -f(x)$，得 $f(x)$ 是奇函数，其图象关于原点对称．又 $f\left(\frac{\pi}{2}\right) = \frac{1 + \frac{\pi}{2}}{\left(\frac{\pi}{2}\right)^2} = \frac{4 + 2\pi}{\pi^2} > 1$, $f(\pi) = \frac{\pi}{-1 + \pi^2} > 0$．故选 $D$．
\end{solution}
\end{question}

\begin{question}[points=5]
我国古代典籍《周易》用“卦”描述万物的变化．每一“重卦”由从下到上排列的 $6$ 个爻组成，爻分为阳爻“——”和阴爻“— —”，如图就是一重卦．在所有重卦中随机取一重卦，则该重卦恰有 $3$ 个阳爻的概率是

\begin{center}
  \includegraphics[width=.38\linewidth]{question-6-1.jpg}
\end{center}
\begin{choices}
  \item $\frac{5}{16}$
  \item $\frac{11}{32}$
  \item $\frac{21}{32}$
  \item $\frac{11}{16}$
\end{choices}
\begin{solution}
\textbf{答案：}$A$

\textbf{分析：}本题主要考查利用两个计数原理与排列组合计算古典概型问题，渗透了传统文化、数学计算等数学素养，“重卦”中每一爻有两种情况，基本事件计算是住店问题，该重卦恰有 $3$ 个阳爻是相同元素的排列问题，利用直接法即可计算．

\textbf{详解：}由题知，每一爻有 $2$ 种情况，一重卦的 $6$ 爻有 $2^6$ 情况，其中 $6$ 爻中恰有 $3$ 个阳爻情况有 $C_6^3$，所以该重卦恰有 $3$ 个阳爻的概率为 $\frac{C_6^3}{2^6} = \frac{5}{16}$，故选 $A$．
\end{solution}
\end{question}

\begin{question}[points=5]
已知非零向量 $\mathbf{a}$, $\mathbf{b}$ 满足 $|\mathbf{a}| = 2|\mathbf{b}|$, 且 $(\mathbf{a} - \mathbf{b}) \perp \mathbf{b}$, 则 $\mathbf{a}$ 与 $\mathbf{b}$ 的夹角为
\begin{choices}
  \item $\frac{\pi}{6}$
  \item $\frac{\pi}{3}$
  \item $\frac{2\pi}{3}$
  \item $\frac{5\pi}{6}$
\end{choices}
\begin{solution}
\textbf{答案：}$B$

\textbf{分析：}本题主要考查利用平面向量数量积计算向量长度、夹角与垂直问题，渗透了转化与化归、数学计算等数学素养．先由 $(\mathbf{a} - \mathbf{b}) \perp \mathbf{b}$ 得出向量 $\mathbf{a}$, $\mathbf{b}$ 的数量积与其模的关系，再利用向量夹角公式即可计算出向量夹角．

\textbf{详解：}因为 $(\mathbf{a} - \mathbf{b}) \perp \mathbf{b}$, 所以 $(\mathbf{a} - \mathbf{b}) \cdot \mathbf{b} = \mathbf{a} \cdot \mathbf{b} - |\mathbf{b}|^2 = 0$, 所以 $\mathbf{a} \cdot \mathbf{b} = |\mathbf{b}|^2$, 所以 $\cos \theta = \frac{\mathbf{a} \cdot \mathbf{b}}{|\mathbf{a}| |\mathbf{b}|} = \frac{|\mathbf{b}|^2}{2|\mathbf{b}|^2} = \frac{1}{2}$, 所以 $\mathbf{a}$ 与 $\mathbf{b}$ 的夹角为 $\frac{\pi}{3}$, 故选 $B$．
\end{solution}
\end{question}

\begin{question}[points=5]
如图是求 $\frac{1}{2 + \frac{1}{2 + \frac{1}{2}}}$ 的程序框图，图中空白框中应填入

\begin{center}
  \includegraphics[width=.38\linewidth]{question-8-1.jpg}
\end{center}
\begin{choices}
  \item $A = \frac{1}{2 + A}$
  \item $A = 2 + \frac{1}{A}$
  \item $A = \frac{1}{1 + 2A}$
  \item $A = 1 + \frac{1}{2A}$
\end{choices}
\begin{solution}
\textbf{答案：}$A$

\textbf{分析：}本题主要考查算法中的程序框图，渗透阅读、分析与解决问题等素养，认真分析式子结构特征与程序框图结构，即可找出作出选择．

\textbf{详解：}执行第 $1$ 次，$A = \frac{1}{2}$, $k = 1 \le 2$ 是，因为第一次应该计算 $\frac{1}{2 + \frac{1}{2}} = \frac{1}{2 + A}$，$k = k + 1 = 2$，循环，执行第 $2$ 次，$k = 2 \le 2$，是，因为第二次应该计算 $\frac{1}{2 + \frac{1}{2 + \frac{1}{2}}} = \frac{1}{2 + A}$，$k = k + 1 = 3$，循环，执行第 $3$ 次，$k = 2 \le 2$，否，输出，故循环体为 $A = \frac{1}{2 + A}$，故选 $A$．
\end{solution}
\end{question}

\begin{question}[points=5]
记 $S_n$ 为等差数列 $\{a_n\}$ 的前 $n$ 项和．已知 $S_4 = 0$, $a_5 = 5$, 则
\begin{choices}
  \item $a_n = 2n - 5$
  \item $a_n = 3n - 10$
  \item $S_n = 2n^2 - 8n$
  \item $S_n = \frac{1}{2}n^2 - 2n$
\end{choices}
\begin{solution}
\textbf{答案：}$A$

\textbf{分析：}等差数列通项公式与前 $n$ 项和公式．本题还可用排除，对 $B$, $a_5 = 5$, $S_4 = \frac{4(-7 + 2)}{2} = -10 \neq 0$，排除 $B$，对 $C$, $S_4 = 0$, $a_5 = S_5 - S_4 = 2 \times 5^2 - 8 \times 5 - 0 = 10 \neq 5$，排除 $C$．对 $D$, $S_4 = 0$, $a_5 = S_5 - S_4 = \frac{1}{2} \times 4^2 - 2 \times 4 - 0 = 0 \neq 5$，排除 $D$，故选 $A$．

\textbf{详解：}由题知，$\begin{cases} S_4 = 4a_1 + \frac{d}{2} \times 4 \times 3 = 0 \\ a_5 = a_1 + 4d = 5 \end{cases}$，解得 $\begin{cases} a_1 = -3 \\ d = 2 \end{cases}$，$\therefore a_n = 2n - 5$，故选 $A$．
\end{solution}
\end{question}

\begin{question}[points=5]
已知椭圆 $C$ 的焦点为 $F_1(-1, 0)$, $F_2(1, 0)$，过 $F_2$ 的直线与 $C$ 交于 $A$, $B$ 两点．若 $|AF_2| = 2|F_2B|$, $|AB| = |BF_1|$, 则 $C$ 的方程为

\begin{center}
  \includegraphics[width=.38\linewidth]{question-10-1.jpg}
\end{center}
\begin{choices}
  \item $\frac{x^2}{2} + y^2 = 1$
  \item $\frac{x^2}{3} + \frac{y^2}{2} = 1$
  \item $\frac{x^2}{4} + \frac{y^2}{3} = 1$
  \item $\frac{x^2}{5} + \frac{y^2}{4} = 1$
\end{choices}
\begin{solution}
\textbf{答案：}$B$

\textbf{分析：}可以运用下面方法求解：如图，由已知可设 $|F_2B| = n$，则 $|AF_2| = 2n$, $|BF_1| = |AB| = 3n$，由椭圆的定义有 $2a = |BF_1| + |BF_2| = 4n$, $\therefore |AF_1| = 2a - |AF_2| = 2n$．在 $\triangle AF_1F_2$ 和 $\triangle BF_1F_2$ 中，由余弦定理得 $\begin{cases} 4n^2 + 4 - 2 \cdot 2n \cdot 2 \cdot \cos \angle AF_2F_1 = 4n^2 \\ n^2 + 4 - 2 \cdot n \cdot 2 \cdot \cos \angle BF_2F_1 = 9n^2 \end{cases}$，又 $\angle AF_2F_1$, $\angle BF_2F_1$ 互补，$\therefore \cos \angle AF_2F_1 + \cos \angle BF_2F_1 = 0$，两式消去 $\cos \angle AF_2F_1$, $\cos \angle BF_2F_1$，得 $3n^2 + 6 = 11n^2$，解得 $n = \frac{\sqrt{3}}{2}$．$\therefore 2a = 4n = 2\sqrt{3}$, $\therefore a = \sqrt{3}$, $\therefore b^2 = a^2 - c^2 = 3 - 1 = 2$, $\therefore$ 所求椭圆方程为 $\frac{x^2}{3} + \frac{y^2}{2} = 1$，故选 $B$．

\textbf{详解：}如图，由已知可设 $|F_2B| = n$，则 $|AF_2| = 2n$, $|BF_1| = |AB| = 3n$，由椭圆的定义有 $2a = |BF_1| + |BF_2| = 4n$, $\therefore |AF_1| = 2a - |AF_2| = 2n$．在 $\triangle AF_1B$ 中，由余弦定理推论得 $\cos \angle F_1AB = \frac{4n^2 + 9n^2 - 9n^2}{2 \cdot 2n \cdot 3n} = \frac{1}{3}$．在 $\triangle AF_1F_2$ 中，由余弦定理得 $4n^2 + 4n^2 - 2 \cdot 2n \cdot 2n \cdot \frac{1}{3} = 4$，解得 $n = \frac{\sqrt{3}}{2}$．$\therefore 2a = 4n = 2\sqrt{3}$, $\therefore a = \sqrt{3}$, $\therefore b^2 = a^2 - c^2 = 3 - 1 = 2$, $\therefore$ 所求椭圆方程为 $\frac{x^2}{3} + \frac{y^2}{2} = 1$，故选 $B$．
\end{solution}
\end{question}

\begin{question}[points=5]
关于函数 $f(x) = \sin|x| + |\sin x|$ 有下述四个结论：

① $f(x)$ 是偶函数

② $f(x)$ 在区间 $\left(\frac{\pi}{2}, \pi\right)$ 单调递增

③ $f(x)$ 在 $[-\pi, \pi]$ 有 $4$ 个零点

④ $f(x)$ 的最大值为 $2$

其中所有正确结论的编号是
\begin{choices}
  \item ①②④
  \item ②④
  \item ①④
  \item ①③
\end{choices}
\begin{solution}
\textbf{答案：}$C$

\textbf{分析：}画出函数 $f(x) = \sin|x| + |\sin x|$ 的图象，由图象可得①④正确，故选 $C$．

\textbf{详解：}$\because f(-x) = \sin|-x| + |\sin(-x)| = \sin|x| + |\sin x| = f(x)$, $\therefore f(x)$ 为偶函数，故①正确．当 $\frac{\pi}{2} < x < \pi$ 时，$f(x) = 2\sin x$，它在区间 $\left(\frac{\pi}{2}, \pi\right)$ 单调递减，故②错误．当 $0 \le x \le \pi$ 时，$f(x) = 2\sin x$，它有两个零点：$0$, $\pi$；当 $-\pi \le x < 0$ 时，$f(x) = \sin(-x) - \sin x = -2\sin x$，它有一个零点：$-\pi$，故 $f(x)$ 在 $[-\pi, \pi]$ 有 $3$ 个零点：$-\pi$, $0$, $\pi$，故③错误．当 $x \in [2k\pi, 2k\pi + \pi] (k \in \mathbb{N}^*)$ 时，$f(x) = 2\sin x$；当 $x \in [2k\pi + \pi, 2k\pi + 2\pi] (k \in \mathbb{N}^*)$ 时，$f(x) = \sin x - \sin x = 0$，又 $f(x)$ 为偶函数，$\therefore f(x)$ 的最大值为 $2$，故④正确．综上所述，①④正确，故选 $C$．
\end{solution}
\end{question}

\begin{question}[points=5]
已知三棱锥 $P-ABC$ 的四个顶点在球 $O$ 的球面上，$PA = PB = PC$，$\triangle ABC$ 是边长为 $2$ 的正三角形，$E$, $F$ 分别是 $PA$, $PB$ 的中点，$\angle CEF = 90^\circ$，则球 $O$ 的体积为
\begin{choices}
  \item $8\sqrt{6}\pi$
  \item $4\sqrt{6}\pi$
  \item $2\sqrt{6}\pi$
  \item $\sqrt{6}\pi$
\end{choices}
\begin{solution}
\textbf{答案：}$D$

\textbf{分析：}本题也可用解三角形方法，达到求出棱长的目的．适合空间想象能力略差学生．设 $PA = PB = PC = 2x$, $E$, $F$ 分别为 $PA$, $AB$ 中点，$\therefore EF \parallel PB$，且 $EF = \frac{1}{2}PB = x$，$\because \triangle ABC$ 为边长为 $2$ 的等边三角形，$\therefore CF = \sqrt{3}$ 又 $\angle CEF = 90^\circ$ $\therefore CE = \sqrt{3 - x^2}$, $AE = \frac{1}{2}PA = x$ $\triangle AEC$ 中余弦定理 $\cos \angle EAC = \frac{x^2 + 4 - (3 - x^2)}{2 \times 2 \times x}$，作 $PD \perp AC$ 于 $D$，$\because PA = PC$，$\therefore D$ 为 $AC$ 中点，$\cos \angle EAC = \frac{AD}{PA} = \frac{1}{2x}$，$\therefore \frac{x^2 + 4 - 3 + x^2}{4x} = \frac{1}{2x}$，$\therefore 2x^2 + 1 = 2$ $\therefore x^2 = \frac{1}{2}$ $x = \frac{\sqrt{2}}{2}$，$\therefore PA = PB = PC = \sqrt{2}$，又 $AB = BC = AC = 2$，$\therefore PA$, $PB$, $PC$ 两两垂直，$\therefore 2R = \sqrt{2 + 2 + 2} = \sqrt{6}$，$\therefore R = \frac{\sqrt{6}}{2}$，$\therefore V = \frac{4}{3}\pi R^3 = \frac{4}{3}\pi \times \frac{6\sqrt{6}}{8} = \sqrt{6}\pi$，故选 $D$．

\textbf{详解：}$\because PA = PB = PC$, $\triangle ABC$ 为边长为 $2$ 的等边三角形，$\therefore P-ABC$ 为正三棱锥，$\therefore PB \perp AC$，又 $E$, $F$ 分别为 $PA$, $AB$ 中点，$\therefore EF \parallel PB$，$\therefore EF \perp AC$，又 $EF \perp CE$, $CE \cap AC = C$, $\therefore EF \perp$ 平面 $PAC$，$PB \perp$ 平面 $PAC$，$\therefore \angle PAB = 90^\circ$, $\therefore PA = PB = PC = \sqrt{2}$，$\therefore P-ABC$ 为正方体一部分，$2R = \sqrt{2 + 2 + 2} = \sqrt{6}$，即 $R = \frac{\sqrt{6}}{2}$, $\therefore V = \frac{4}{3}\pi R^3 = \frac{4}{3}\pi \times \frac{6\sqrt{6}}{8} = \sqrt{6}\pi$，故选 $D$．
\end{solution}
\end{question}

\section*{二、填空题}
本题共4小题，每小题5分，共20分。

\begin{question}[points=5]
曲线 $y = 3(x^2 + x)e^x$ 在点 $(0, 0)$ 处的切线方程为．
\fillin[{$3x - y = 0$}]
\begin{solution}
\textbf{答案：}$3x - y = 0$

\textbf{分析：}本题根据导数的几何意义，通过求导数，确定得到切线的斜率，利用直线方程的点斜式求得切线方程．

\textbf{详解：}$y' = 3(2x + 1)e^x + 3(x^2 + x)e^x = 3(x^2 + 3x + 1)e^x$，所以，$k = y'|_{x=0} = 3$，所以，曲线 $y = 3(x^2 + x)e^x$ 在点 $(0, 0)$ 处的切线方程为 $y = 3x$，即 $3x - y = 0$．
\end{solution}
\end{question}

\begin{question}[points=5]
记 $S_n$ 为等比数列 $\{a_n\}$ 的前 $n$ 项和．若 $a_1 = \frac{1}{3}$, $a_4^2 = a_6$, 则 $S_5 =$ ．
\fillin[{$\frac{121}{3}$}]
\begin{solution}
\textbf{答案：}$\frac{121}{3}$

\textbf{分析：}本题根据已知条件，列出关于等比数列公比 $q$ 的方程，应用等比数列的求和公式，计算得到 $S_5$．题目的难度不大，注重了基础知识、基本计算能力的考查．

\textbf{详解：}设等比数列的公比为 $q$，由已知 $a_1 = \frac{1}{3}$, $a_4^2 = a_6$, 所以 $\left(\frac{1}{3}q^3\right)^2 = \frac{1}{3}q^5$, 又 $q \neq 0$, 所以 $q = 3$, 所以 $S_5 = \frac{a_1(1 - q^5)}{1 - q} = \frac{\frac{1}{3}(1 - 3^5)}{1 - 3} = \frac{121}{3}$．
\end{solution}
\end{question}

\begin{question}[points=5]
甲、乙两队进行篮球决赛，采取七场四胜制（当一队赢得四场胜利时，该队获胜，决赛结束）．根据前期比赛成绩，甲队的主客场安排依次为“主主客客主客主”．设甲队主场取胜的概率为 $0.6$，客场取胜的概率为 $0.5$，且各场比赛结果相互独立，则甲队以 $4 \colon 1$ 获胜的概率是．
\fillin[{$0.216$}]
\begin{solution}
\textbf{答案：}$0.216$

\textbf{分析：}本题应注意分情况讨论，即前五场甲队获胜的两种情况，应用独立事件的概率的计算公式求解．题目有一定的难度，注重了基础知识、基本计算能力及分类讨论思想的考查．

\textbf{详解：}前五场中有一场客场输时，甲队以 $4 \colon 1$ 获胜的概率是 $0.6^3 \times 0.5 \times 0.5 \times 2 = 0.108$, 前五场中有一场主场输时，甲队以 $4 \colon 1$ 获胜的概率是 $0.4 \times 0.6^2 \times 0.5^2 \times 3 = 0.108$, 综上所述，甲队以 $4 \colon 1$ 获胜的概率是 $0.108 + 0.108 = 0.216$．
\end{solution}
\end{question}

\begin{question}[points=5]
已知双曲线 $C$：$\frac{x^2}{a^2} - \frac{y^2}{b^2} = 1$ $(a > 0, b > 0)$ 的左、右焦点分别为 $F_1$, $F_2$，过 $F_1$ 的直线与 $C$ 的两条渐近线分别交于 $A$, $B$ 两点．若 $\overrightarrow{F_1A} = \overrightarrow{AB}$, $\overrightarrow{F_1B} \cdot \overrightarrow{F_2B} = 0$，则 $C$ 的离心率为．
\fillin[{$2$}]
\begin{solution}
\textbf{答案：}$2$

\textbf{分析：}本题考查平面向量结合双曲线的渐进线和离心率，渗透了逻辑推理、直观想象和数学运算素养．采取几何法，利用数形结合思想解题．

\textbf{详解：}如图，由 $\overrightarrow{F_1A} = \overrightarrow{AB}$, 得 $|F_1A| = |AB|$. 又 $|OF_1| = |OF_2|$, 得 $OA$ 是三角形 $F_1F_2B$ 的中位线，即 $BF_2 \parallel OA$, $|BF_2| = 2|OA|$. 由 $\overrightarrow{F_1B} \cdot \overrightarrow{F_2B} = 0$, 得 $F_1B \perp F_2B$, $OA \perp F_1A$, 则 $|OB| = |OF_1| = |OF_2|$, 有 $\angle OBF_2 = \angle BF_2O = 2\angle OBF_1 = 2\angle OF_1B$, $\angle AOB = \angle AOF_1$. 又 $OA$ 与 $OB$ 都是渐近线，得 $\angle BOF_2 = \angle AOF_1$, 则 $\angle BOF_2 = 60^\circ$. 又渐近线 $OB$ 的斜率为 $\frac{b}{a} = \tan 60^\circ = \sqrt{3}$, 所以该双曲线的离心率为 $e = \frac{c}{a} = \sqrt{1 + \left(\frac{b}{a}\right)^2} = \sqrt{1 + (\sqrt{3})^2} = 2$．
\end{solution}
\end{question}

\section*{三、解答题}
共70分。解答应写出文字说明、证明过程或演算步骤。第17~21题为必考题，每个试题考生都必须作答。第22、23题为选考题，考生根据要求作答。

\begin{question}[points=12]
$\triangle ABC$ 的内角 $A$, $B$, $C$ 的对边分别为 $a$, $b$, $c$, 设 $(\sin B - \sin C)^2 = \sin^2 A - \sin B \sin C$．

(1) 求 $A$;

(2) 若 $\sqrt{2}a + b = 2c$, 求 $\sin C$．
\begin{solution}
\textbf{答案：}(1) $A = \frac{\pi}{3}$; (2) $\sin C = \frac{\sqrt{6} + \sqrt{2}}{4}$

\textbf{分析：}(1) 利用正弦定理化简已知边角关系式可得：$b^2 + c^2 - a^2 = bc$，从而可整理出 $\cos A$，根据 $A \in (0, \pi)$ 可求得结果；(2) 利用正弦定理可得 $\sqrt{2}\sin A + \sin B = 2\sin C$，利用 $\sin B = \sin(A + C)$、两角和差正弦公式可得关于 $\sin C$ 和 $\cos C$ 的方程，结合同角三角函数关系解方程可求得结果．

\textbf{详解：}(1) $(\sin B - \sin C)^2 = \sin^2 B - 2\sin B\sin C + \sin^2 C = \sin^2 A - \sin B\sin C$，即 $\sin^2 B + \sin^2 C - \sin^2 A = \sin B\sin C$，由正弦定理可得：$b^2 + c^2 - a^2 = bc$，$\therefore \cos A = \frac{b^2 + c^2 - a^2}{2bc} = \frac{1}{2}$，$\because A \in (0, \pi)$，$\therefore A = \frac{\pi}{3}$．
(2) $\because \sqrt{2}a + b = 2c$，由正弦定理得：$\sqrt{2}\sin A + \sin B = 2\sin C$，又 $\sin B = \sin(A + C) = \sin A \cos C + \cos A \sin C$，$A = \frac{\pi}{3}$，$\therefore \sqrt{2} \times \frac{\sqrt{3}}{2} + \frac{\sqrt{3}}{2}\cos C + \frac{1}{2}\sin C = 2\sin C$，整理可得：$3\sin C - \sqrt{6} = 3\cos C$，$\because \sin^2 C + \cos^2 C = 1$，$\therefore (3\sin C - \sqrt{6})^2 = 9(1 - \sin^2 C)$，解得：$\sin C = \frac{\sqrt{6} + \sqrt{2}}{4}$ 或 $\frac{\sqrt{6} - \sqrt{2}}{4}$．因为 $\sin B = 2\sin C - \sqrt{2}\sin A = 2\sin C - \frac{\sqrt{6}}{2} > 0$，所以 $\sin C > \frac{\sqrt{6}}{4}$，故 $\sin C = \frac{\sqrt{6} + \sqrt{2}}{4}$．
\end{solution}
\end{question}

\begin{question}[points=12]
如图，直四棱柱 $ABCD-A_1B_1C_1D_1$ 的底面是菱形，$AA_1 = 4$, $AB = 2$, $\angle BAD = 60^\circ$, $E$, $M$, $N$ 分别是 $BC$, $BB_1$, $A_1D$ 的中点．

(1) 证明：$MN \parallel$ 平面 $C_1DE$;

(2) 求二面角 $A-MA_1-N$ 的正弦值．

\begin{center}
  \includegraphics[width=.38\linewidth]{question-18-1.jpg}
\end{center}
\begin{solution}
\textbf{答案：}(1) 见解析; (2) $\frac{\sqrt{10}}{5}$

\textbf{分析：}(1) 利用三角形中位线和 $A_1D \parallel B_1C$ 可证得 $ME \parallel ND$，证得四边形 $MNDE$ 为平行四边形，进而证得 $MN \parallel DE$，根据线面平行判定定理可证得结论；(2) 以菱形 $ABCD$ 对角线交点为原点可建立空间直角坐标系，通过取 $AB$ 中点 $F$，可证得 $DF \perp$ 平面 $AMA_1$，得到平面 $AMA_1$ 的法向量 $\overrightarrow{DF}$；再通过向量法求得平面 $MA_1N$ 的法向量 $\mathbf{n}$，利用向量夹角公式求得两个法向量夹角的余弦值，进而可求得所求二面角的正弦值．

\textbf{详解：}(1) 连接 $ME$, $B_1C$．$\because M$, $E$ 分别为 $BB_1$, $BC$ 中点，$\therefore ME$ 为 $\triangle B_1BC$ 的中位线，$\therefore ME \parallel B_1C$ 且 $ME = \frac{1}{2}B_1C$．又 $N$ 为 $A_1D$ 中点，且 $A_1D \parallel B_1C$，$\therefore ND \parallel B_1C$ 且 $ND = \frac{1}{2}B_1C$，$\therefore ME \parallel ND$，$\therefore$ 四边形 $MNDE$ 为平行四边形，$\therefore MN \parallel DE$，又 $MN \not\subset$ 平面 $C_1DE$, $DE \subset$ 平面 $C_1DE$，$\therefore MN \parallel$ 平面 $C_1DE$．
(2) 设 $AC \cap BD = O$, $A_1C_1 \cap B_1D_1 = O_1$，由直四棱柱性质可知：$OO_1 \perp$ 平面 $ABCD$．$\because$ 四边形 $ABCD$ 为菱形，$\therefore AC \perp BD$．则以 $O$ 为原点，可建立如图所示的空间直角坐标系．则 $A(\sqrt{3}, 0, 0)$, $M(0, 1, 2)$, $A_1(\sqrt{3}, 0, 4)$, $D(0, -1, 0)$, $N\left(\frac{\sqrt{3}}{2}, -\frac{1}{2}, 2\right)$．取 $AB$ 中点 $F$，连接 $DF$，则 $F\left(\frac{\sqrt{3}}{2}, \frac{1}{2}, 0\right)$．$\because$ 四边形 $ABCD$ 为菱形且 $\angle BAD = 60^\circ$，$\therefore \triangle BAD$ 为等边三角形，$\therefore DF \perp AB$．又 $AA_1 \perp$ 平面 $ABCD$, $DF \subset$ 平面 $ABCD$，$\therefore DF \perp AA_1$，$\therefore DF \perp$ 平面 $ABB_1A_1$，即 $DF \perp$ 平面 $AMA_1$，$\therefore \overrightarrow{DF}$ 为平面 $AMA_1$ 的一个法向量，且 $\overrightarrow{DF} = \left(\frac{\sqrt{3}}{2}, \frac{3}{2}, 0\right)$．设平面 $MA_1N$ 的法向量 $\mathbf{n} = (x, y, z)$，又 $\overrightarrow{MA_1} = (\sqrt{3}, -1, 2)$, $\overrightarrow{MN} = \left(\frac{\sqrt{3}}{2}, -\frac{3}{2}, 0\right)$，$\because \begin{cases} \mathbf{n} \cdot \overrightarrow{MA_1} = \sqrt{3}x - y + 2z = 0 \\ \mathbf{n} \cdot \overrightarrow{MN} = \frac{\sqrt{3}}{2}x - \frac{3}{2}y = 0 \end{cases}$，令 $x = \sqrt{3}$，则 $y = 1$, $z = -1$，$\therefore \mathbf{n} = (\sqrt{3}, 1, -1)$．$\therefore \cos \langle \overrightarrow{DF}, \mathbf{n} \rangle = \frac{\overrightarrow{DF} \cdot \mathbf{n}}{|\overrightarrow{DF}| |\mathbf{n}|} = \frac{3}{\sqrt{15} \times \sqrt{5}} = \frac{\sqrt{3}}{5}$，$\therefore \sin \langle \overrightarrow{DF}, \mathbf{n} \rangle = \frac{\sqrt{10}}{5}$．$\therefore$ 二面角 $A-MA_1-N$ 的正弦值为 $\frac{\sqrt{10}}{5}$．
\end{solution}
\end{question}

\begin{question}[points=12]
已知抛物线 $C$：$y^2 = 3x$ 的焦点为 $F$，斜率为 $\frac{3}{2}$ 的直线 $l$ 与 $C$ 的交点为 $A$, $B$，与 $x$ 轴的交点为 $P$．

(1) 若 $|AF| + |BF| = 4$，求 $l$ 的方程；

(2) 若 $\overrightarrow{AP} = 3\overrightarrow{PB}$，求 $|AB|$．
\begin{solution}
\textbf{答案：}(1) $12x - 8y - 7 = 0$; (2) $\frac{4\sqrt{13}}{3}$

\textbf{分析：}(1) 设直线 $l$：$y = \frac{3}{2}x + m$，$A(x_1, y_1)$, $B(x_2, y_2)$；根据抛物线焦半径公式可得 $x_1 + x_2 = 1$；联立直线方程与抛物线方程，利用韦达定理可构造关于 $m$ 的方程，解方程求得结果；(2) 设直线 $l$：$x = \frac{2}{3}y + t$；联立直线方程与抛物线方程，得到韦达定理的形式；利用 $\overrightarrow{AP} = 3\overrightarrow{PB}$ 可得 $y_1 = -3y_2$，结合韦达定理可求得 $y_1y_2$；根据弦长公式可求得结果．

\textbf{详解：}(1) 设直线 $l$ 方程为：$y = \frac{3}{2}x + m$, $A(x_1, y_1)$, $B(x_2, y_2)$．由抛物线焦半径公式可知：$|AF| + |BF| = x_1 + x_2 + \frac{3}{2} = 4$，$\therefore x_1 + x_2 = \frac{5}{2}$．联立 $\begin{cases} y = \frac{3}{2}x + m \\ y^2 = 3x \end{cases}$ 得：$9x^2 + (12m - 12)x + 4m^2 = 0$，则 $\Delta = (12m - 12)^2 - 144m^2 > 0$，$\therefore m < \frac{1}{2}$，$\therefore x_1 + x_2 = -\frac{12m - 12}{9} = \frac{5}{2}$，解得：$m = -\frac{7}{8}$．$\therefore$ 直线 $l$ 的方程为：$y = \frac{3}{2}x - \frac{7}{8}$，即 $12x - 8y - 7 = 0$．
(2) 设 $P(t, 0)$，则可设直线 $l$ 方程为：$x = \frac{2}{3}y + t$．联立 $\begin{cases} x = \frac{2}{3}y + t \\ y^2 = 3x \end{cases}$ 得：$y^2 - 2y - 3t = 0$，则 $\Delta = 4 + 12t > 0$，$\therefore t > -\frac{1}{3}$，$\therefore y_1 + y_2 = 2$, $y_1y_2 = -3t$．$\because \overrightarrow{AP} = 3\overrightarrow{PB}$，$\therefore y_1 = -3y_2$，$\therefore y_2 = -1$, $y_1 = 3$，$\therefore y_1y_2 = -3$．则 $|AB| = \sqrt{1 + \frac{4}{9}} \cdot \sqrt{(y_1 + y_2)^2 - 4y_1y_2} = \frac{\sqrt{13}}{3} \cdot \sqrt{4 + 12} = \frac{4\sqrt{13}}{3}$．
\end{solution}
\end{question}

\begin{question}[points=12]
已知函数 $f(x) = \sin x - \ln(1 + x)$, $f'(x)$ 为 $f(x)$ 的导数．证明：

(1) $f'(x)$ 在区间 $\left(-1, \frac{\pi}{2}\right)$ 存在唯一极大值点；

(2) $f(x)$ 有且仅有 $2$ 个零点．
\begin{solution}
\textbf{答案：}见解析

\textbf{分析：}(1) 求得导函数后，可判断出导函数在 $\left(-1, \frac{\pi}{2}\right)$ 上单调递减，根据零点存在定理可判断出 $\exists x_0 \in \left(0, \frac{\pi}{2}\right)$，使得 $g'(x_0) = 0$，进而得到导函数在 $\left(-1, \frac{\pi}{2}\right)$ 上的单调性，从而可证得结论；(2) 由（1）的结论可知 $x = 0$ 为 $f(x)$ 在 $(-1, 0]$ 上的唯一零点；当 $x \in \left(0, \frac{\pi}{2}\right]$ 时，首先可判断出在 $(0, x_0)$ 上无零点，再利用零点存在定理得到 $f(x)$ 在 $\left(x_0, \frac{\pi}{2}\right)$ 上的单调性，可知 $f(x) > 0$，不存在零点；当 $x \in \left[\frac{\pi}{2}, \pi\right]$ 时，利用零点存在定理和 $f(x)$ 单调性可判断出存在唯一一个零点；当 $x \in (\pi, +\infty)$，可证得 $f(x) < 0$；综合上述情况可证得结论．

\textbf{详解：}(1) 由题意知：$f(x)$ 定义域为 $(-1, +\infty)$ 且 $f'(x) = \cos x - \frac{1}{x+1}$．令 $g(x) = \cos x - \frac{1}{x+1}$, $x \in \left(-1, \frac{\pi}{2}\right)$，$\therefore g'(x) = -\sin x + \frac{1}{(x+1)^2}$, $x \in \left(-1, \frac{\pi}{2}\right)$．$\because \frac{1}{(x+1)^2}$ 在 $\left(-1, \frac{\pi}{2}\right)$ 上单调递减，$-\sin x$ 在 $\left(-1, \frac{\pi}{2}\right)$ 上单调递减，$\therefore g'(x)$ 在 $\left(-1, \frac{\pi}{2}\right)$ 上单调递减．又 $g'(0) = -\sin 0 + 1 = 1 > 0$, $g'\left(\frac{\pi}{2}\right) = -\sin\frac{\pi}{2} + \frac{4}{(\pi+2)^2} = -1 + \frac{4}{(\pi+2)^2} < 0$，$\therefore \exists x_0 \in \left(0, \frac{\pi}{2}\right)$，使得 $g'(x_0) = 0$．$\therefore$ 当 $x \in (-1, x_0)$ 时，$g'(x) > 0$；当 $x \in \left(x_0, \frac{\pi}{2}\right)$ 时，$g'(x) < 0$．即 $g(x)$ 在 $(-1, x_0)$ 上单调递增；在 $\left(x_0, \frac{\pi}{2}\right)$ 上单调递减，则 $x = x_0$ 为 $g(x)$ 唯一的极大值点，即 $f'(x)$ 在区间 $\left(-1, \frac{\pi}{2}\right)$ 上存在唯一的极大值点 $x_0$．
(2) 由（1）知：$f'(x) = \cos x - \frac{1}{x+1}$, $x \in (-1, +\infty)$．
①当 $x \in (-1, 0]$ 时，由（1）可知 $f'(x)$ 在 $(-1, 0]$ 上单调递增，$\therefore f'(x) \le f'(0) = 0$，$\therefore f(x)$ 在 $(-1, 0]$ 上单调递减，又 $f(0) = 0$，$\therefore x = 0$ 为 $f(x)$ 在 $(-1, 0]$ 上的唯一零点．
②当 $x \in \left(0, \frac{\pi}{2}\right]$ 时，$f'(x)$ 在 $(0, x_0)$ 上单调递增，在 $\left(x_0, \frac{\pi}{2}\right)$ 上单调递减，又 $f'(0) = 0$，$\therefore f'(x_0) > 0$，$\therefore f(x)$ 在 $(0, x_0)$ 上单调递增，此时 $f(x) > f(0) = 0$，不存在零点．又 $f'\left(\frac{\pi}{2}\right) = \cos\frac{\pi}{2} - \frac{2}{\pi+2} = -\frac{2}{\pi+2} < 0$，$\therefore \exists x_1 \in \left(x_0, \frac{\pi}{2}\right)$，使得 $f'(x_1) = 0$，$\therefore f(x)$ 在 $(x_0, x_1)$ 上单调递增，在 $\left(x_1, \frac{\pi}{2}\right)$ 上单调递减．又 $f(x_0) > f(0) = 0$, $f\left(\frac{\pi}{2}\right) = \sin\frac{\pi}{2} - \ln\left(1 + \frac{\pi}{2}\right) = \ln\frac{2e}{\pi+2} > \ln 1 = 0$，$\therefore f(x) > 0$ 在 $\left(x_0, \frac{\pi}{2}\right)$ 上恒成立，此时不存在零点．
③当 $x \in \left[\frac{\pi}{2}, \pi\right]$ 时，$\sin x$ 单调递减，$-\ln(x+1)$ 单调递减，$\therefore f(x)$ 在 $\left[\frac{\pi}{2}, \pi\right]$ 上单调递减．又 $f\left(\frac{\pi}{2}\right) > 0$, $f(\pi) = \sin\pi - \ln(\pi+1) = -\ln(\pi+1) < 0$，即 $f(\pi) \cdot f\left(\frac{\pi}{2}\right) < 0$，又 $f(x)$ 在 $\left[\frac{\pi}{2}, \pi\right]$ 上单调递减，$\therefore f(x)$ 在 $\left[\frac{\pi}{2}, \pi\right]$ 上存在唯一零点．
④当 $x \in (\pi, +\infty)$ 时，$\sin x \in [-1, 1]$, $\ln(x+1) > \ln(\pi+1) > \ln e = 1$，$\therefore \sin x - \ln(x+1) < 0$，即 $f(x)$ 在 $(\pi, +\infty)$ 上不存在零点．
综上所述：$f(x)$ 有且仅有 $2$ 个零点．
\end{solution}
\end{question}

\begin{question}[points=12]
为了治疗某种疾病，研制了甲、乙两种新药，希望知道哪种新药更有效，为此进行动物试验．试验方案如下：每一轮选取两只白鼠对药效进行对比试验．对于两只白鼠，随机选一只施以甲药，另一只施以乙药．一轮的治疗结果得出后，再安排下一轮试验．当其中一种药治愈的白鼠比另一种药治愈的白鼠多 $4$ 只时，就停止试验，并认为治愈只数多的药更有效．为了方便描述问题，约定：对于每轮试验，若施以甲药的白鼠治愈且施以乙药的白鼠未治愈则甲药得 $1$ 分，乙药得 $-1$ 分；若施以乙药的白鼠治愈且施以甲药的白鼠未治愈则乙药得 $1$ 分，甲药得 $-1$ 分；若都治愈或都未治愈则两种药均得 $0$ 分．甲、乙两种药的治愈率分别记为 $\alpha$ 和 $\beta$，一轮试验中甲药的得分记为 $X$．

(1) 求 $X$ 的分布列；

(2) 若甲药、乙药在试验开始时都赋予 $4$ 分，$p_i$ $(i = 0, 1, \cdots, 8)$ 表示“甲药的累计得分为 $i$ 时，最终认为甲药比乙药更有效”的概率，则 $p_0 = 0$, $p_8 = 1$, $p_i = a p_{i-1} + b p_i + c p_{i+1}$ $(i = 1, 2, \cdots, 7)$，其中 $a = P(X = -1)$, $b = P(X = 0)$, $c = P(X = 1)$. 假设 $\alpha = 0.5$, $\beta = 0.8$.

(i) 证明：$\{p_{i+1} - p_i\}$ $(i = 0, 1, 2, \cdots, 7)$ 为等比数列；

(ii) 求 $p_4$，并根据 $p_4$ 的值解释这种试验方案的合理性．
\begin{solution}
\textbf{答案：}(1) 见解析; (2)(i) 见解析; (ii) $p_4 = \frac{1}{257}$

\textbf{分析：}(1) 首先确定 $X$ 所有可能的取值，再来计算出每个取值对应的概率，从而可得分布列；(2)(i) 求解出 $a$, $b$, $c$ 的取值，可得 $p_i = 0.4p_{i-1} + 0.5p_i + 0.1p_{i+1}$ $(i = 1,2,\cdots,7)$，从而整理出符合等比数列定义的形式，问题得证；(ii) 列出证得的等比数列的通项公式，采用累加的方式，结合 $p_8$ 和 $p_0$ 的值可求得 $p_1$；再次利用累加法可求出 $p_4$．

\textbf{详解：}(1) 由题意可知 $X$ 所有可能的取值为：$-1$, $0$, $1$．$P(X = -1) = (1 - \alpha)\beta$; $P(X = 0) = \alpha\beta + (1 - \alpha)(1 - \beta)$; $P(X = 1) = \alpha(1 - \beta)$．则 $X$ 的分布列如下：
\begin{center}
\begin{tabular}{c|c|c|c}
$X$ & $-1$ & $0$ & $1$ \\ \hline
$P$ & $(1-\alpha)\beta$ & $\alpha\beta + (1-\alpha)(1-\beta)$ & $\alpha(1-\beta)$
\end{tabular}
\end{center}
(2) $\because \alpha = 0.5$, $\beta = 0.8$, $\therefore a = 0.5 \times 0.8 = 0.4$, $b = 0.5 \times 0.8 + 0.5 \times 0.2 = 0.5$, $c = 0.5 \times 0.2 = 0.1$．
(i) $\because p_i = a p_{i-1} + b p_i + c p_{i+1}$ $(i = 1,2,\cdots,7)$，即 $p_i = 0.4p_{i-1} + 0.5p_i + 0.1p_{i+1}$, 整理可得：$5p_i = 4p_{i-1} + p_{i+1}$, $\therefore p_{i+1} - p_i = 4(p_i - p_{i-1})$ $(i = 1,2,\cdots,7)$．又 $p_1 - p_0 = p_1 \neq 0$, $\therefore \{p_{i+1} - p_i\}$ $(i = 0,1,2,\cdots,7)$ 为公比为 $4$，首项为 $p_1$ 的等比数列．
(ii) 由 (i) 可得 $p_8 = p_8 - p_7 + p_7 - p_6 + \cdots + p_1 - p_0 + p_0 = (p_8 - p_7) + (p_7 - p_6) + \cdots + (p_1 - p_0) = \frac{4^8 - 1}{3}p_1$．由于 $p_8 = 1$, 故 $p_1 = \frac{3}{4^8 - 1}$, 所以 $p_4 = (p_4 - p_3) + (p_3 - p_2) + (p_2 - p_1) + (p_1 - p_0) = \frac{4^4 - 1}{3}p_1 = \frac{1}{257}$．$p_4$ 表示最终认为甲药更有效的概率，由计算结果可以看出，在甲药治愈率为 $0.5$，乙药治愈率为 $0.8$ 时，认为甲药更有效的概率为 $p_4 = \frac{1}{257} \approx 0.0039$，此时得出错误结论的概率非常小，说明这种试验方案合理．
\end{solution}
\end{question}

\section*{（二）选考题}
共10分。请考生在第22、23题中任选一题作答。如果多做，则按所做的第一题计分。

\begin{question}[points=10]
[选修 4—4：坐标系与参数方程]

在直角坐标系 $xOy$ 中，曲线 $C$ 的参数方程为 $\begin{cases} x = \dfrac{1 - t^2}{1 + t^2}, \\ y = \dfrac{4t}{1 + t^2} \end{cases}$ ($t$ 为参数)．以坐标原点 $O$ 为极点，$x$ 轴的正半轴为极轴建立极坐标系，直线 $l$ 的极坐标方程为 $2\rho\cos\theta + \sqrt{3}\rho\sin\theta + 11 = 0$．

(1) 求 $C$ 和 $l$ 的直角坐标方程；

(2) 求 $C$ 上的点到 $l$ 距离的最小值．
\begin{solution}
\textbf{答案：}(1) $C: x^2 + \frac{y^2}{4} = 1$ ($x \neq -1$); $l: 2x + \sqrt{3}y + 11 = 0$; (2) $\sqrt{7}$

\textbf{分析：}(1) 利用代入消元法，可求得 $C$ 的直角坐标方程；根据极坐标与直角坐标互化原则可得 $l$ 的直角坐标方程；(2) 利用参数方程表示出 $C$ 上点的坐标，根据点到直线距离公式可将所求距离表示为三角函数的形式，从而根据三角函数的范围可求得最值．

\textbf{详解：}(1) 由 $x = \frac{1 - t^2}{1 + t^2}$ 得：$t^2 = \frac{1 - x}{1 + x}$，又 $y = \frac{4t}{1 + t^2}$, $\therefore y^2 = \frac{16t^2}{(1 + t^2)^2} = \frac{16 \cdot \frac{1-x}{1+x}}{\left(1 + \frac{1-x}{1+x}\right)^2} = 4(1+x)(1-x) = 4 - 4x^2$, 整理可得 $C$ 的直角坐标方程为 $x^2 + \frac{y^2}{4} = 1$ ($x \neq -1$)．又 $x = \rho\cos\theta$, $y = \rho\sin\theta$, $\therefore l$ 的直角坐标方程为 $2x + \sqrt{3}y + 11 = 0$．
(2) 设 $C$ 上点的坐标为 $(\cos\theta, 2\sin\theta)$，则 $C$ 上的点到直线 $l$ 的距离 $d = \frac{|2\cos\theta + 2\sqrt{3}\sin\theta + 11|}{\sqrt{7}} = \frac{4\sin\left(\theta + \frac{\pi}{6}\right) + 11}{\sqrt{7}}$．当 $\sin\left(\theta + \frac{\pi}{6}\right) = -1$ 时，$d$ 取最小值 $\frac{-4 + 11}{\sqrt{7}} = \sqrt{7}$．
\end{solution}
\end{question}

\begin{question}[points=10]
[选修 4—5：不等式选讲]

已知 $a$, $b$, $c$ 为正数，且满足 $abc = 1$．证明：

(1) $\frac{1}{a} + \frac{1}{b} + \frac{1}{c} \le a^2 + b^2 + c^2$;

(2) $(a + b)^3 + (b + c)^3 + (c + a)^3 \ge 24$．
\begin{solution}
\textbf{答案：}见解析

\textbf{分析：}(1) 利用 $abc = 1$ 将所证不等式可变为证明：$a^2 + b^2 + c^2 \ge bc + ac + ab$，利用基本不等式可证得 $2(a^2 + b^2 + c^2) \ge 2ab + 2bc + 2ac$，从而得到结论；(2) 利用基本不等式可得 $(a+b)^3 + (b+c)^3 + (c+a)^3 \ge 3(a+b)(b+c)(c+a)$，再次利用基本不等式可将式转化为 $(a+b)^3 + (b+c)^3 + (c+a)^3 \ge 24(abc)^2$，在取等条件一致的情况下，可得结论．

\textbf{详解：}(1) $\because abc = 1$, $\therefore \frac{1}{a} + \frac{1}{b} + \frac{1}{c} = \left(\frac{1}{a} + \frac{1}{b} + \frac{1}{c}\right) \cdot abc = bc + ac + ab$．$\because 2(a^2 + b^2 + c^2) = (a^2 + b^2) + (b^2 + c^2) + (c^2 + a^2) \ge 2ab + 2bc + 2ac$, 当且仅当 $a = b = c$ 时取等号，$\therefore 2(a^2 + b^2 + c^2) \ge 2(ab + bc + ac)$, 即 $a^2 + b^2 + c^2 \ge ab + bc + ac = \frac{1}{a} + \frac{1}{b} + \frac{1}{c}$．
(2) $\because (a+b)^3 + (b+c)^3 + (c+a)^3 \ge 3(a+b)(b+c)(c+a)$, 当且仅当 $a = b = c$ 时取等号．又 $a+b \ge 2\sqrt{ab}$, $b+c \ge 2\sqrt{bc}$, $a+c \ge 2\sqrt{ac}$ (当且仅当 $a = b = c$ 时等号同时成立)，$\therefore (a+b)^3 + (b+c)^3 + (c+a)^3 \ge 3 \times 2\sqrt{ab} \times 2\sqrt{bc} \times 2\sqrt{ac} = 24(abc)^2 = 24$．
\end{solution}
\end{question}

